
\begin{frame}{Qu'est-ce que Spark?}

Un \alert{moteur d'exécution} basé sur des opérateurs de haut niveau, comme Pig.

\vfill

Comprend des opérateurs Map/Reduce, et d'autres opérateurs de second ordre.

\vfill

\alert{Introduit un concept de collection résidente en mémoire} (RDD) 
qui améliore considérablement certains traitements, dont ceux basés sur des itérations.

\vfill
\begin{blocgauche}
De nombreuses librairies pour la fouille de données (MLib), le traitement des graphes,
le traitement de flux (\textit{streaming}).
\end{blocgauche}
\end{frame}

\section{Les RDD}

\begin{frame}{Les \textit{Resilient Distributed Datasets} (RDD)}

C'est le concept central : \alert{Un RDD est une collection (pour en rester à notre vocabulaire) calculée
à partir d'une source de données} (MongoDB, un flux, un autre RDD) .

\vfill

Un RDD peut être marqué comme \alert{persistant}: il est alors placé en mémoire RAM et conservé par Spark.
\vfill

Spark conserve
\alert{l'historique des opérations} qui a permis de constituer un RDD, et la reprise sur panne s'appuie
sur  cet historique afin de reconstituer  le RDD en cas de panne.

\vfill
\begin{blocgauche}
Un RDD est un "bloc" \alert{non modifiable}. Si nécessaire il est \alert{entièrement recaculé}.
\end{blocgauche}
\end{frame}


\begin{frame}{Un \textit{workflow} avec RDD dans Spark}

Des \alert{transformations} (opérateurs comme dans Pig) créent des RDD à partir d'une ou deux sources
de données.

\vfill

\figSlideWithSize{spark-rdd}{10cm}

\vfill
\begin{blocgauche}
Les RDD persistants sont en préservés en mémoire RAM, et peuvent être réutilisés par plusieurs 
traitements.
\end{blocgauche}
\end{frame}

\begin{frame}{Exemple: analyse de fichiers log}
 
On veut analyser le fichier journal (*log*) d'une application dont un des modules (M) est suspect.

\vfill

On construit  un programme qui charge le  *log*, ne conserve
que les messsages produits par le module *M* et les analyse.

\vfill

On peut analyser par produit, par utilisateur, par période, etc.


\vfill

\figSlideWithSize{spark-log}{9cm}

\end{frame}


\begin{frame}[fragile]{Spécification avec Spark}

Première phase pour construire \texttt{logM}

\begin{minted}{scala}
    // Chargement de la collection
    log = load ("app.log") as (...)
    // Filtrage des messages du module M
    logM = filter log with log.message.contains ("M")    
    // On rend logM persistant !
    logM.persist();
\end{minted}

\vfill

Analyse à partir de \texttt{logM}

\begin{minted}{scala}
   // Filtrage par produit
    logProduit = filter logM 
       with log.message.contains ("product P")   
    // .. analyse du contenu de logProduit     
\end{minted}

\end{frame}

\section{Reprise sur panne}

\begin{frame}{Reprise sur panne dans Spark}

Un RDD est une collection \alert{partitionnée}.

\vfill

\figSlideWithSize{spark-failover}{8cm}

\vfill

\begin{blocgauche}
Panne affecte un calcul basé sur un fragment $F$  persistant:  on relance à partir de $F$.

\medskip

Panne  affecte le fragment $F$  persistant, Spark re-calcule $F$
\end{blocgauche}
\end{frame}

\section{Une session introductive}

\begin{frame}[fragile]{Installation}

Très simple: récupérez la version courante à \url{http://spark.apache.org}.

\vfill

Décompressez dans un répertoire \texttt{spark}.

\vfill

Lancez l'intepréteur.


\begin{minted}{bash}
./bin/spark-shell

scala>
\end{minted}

\vfill

C'est un interpréteur Scala: apprenez le langage Scala! 

\vfill
 
On va implanter le compteur de termes: récupérer un fichier texte quelconque (ou \texttt{loups.txt}
sur le site.)
\end{frame}


\begin{frame}[fragile]{Commandes de base}

Chargement d'un fichier texte dans un RDD (un document par ligne).

\begin{minted}{scala}
val loupsEtMoutons = sc.textFile("loups.txt")
\end{minted}

\vfill

Quelques \alert{actions} appliquées au RDD.

\begin{minted}{scala}
loupsEtMoutons.count() // Nombre de documents dans ce RDD
loupsEtMoutons.first() // Le premier document
loupsEtMoutons.collect() // Tous les documents du RDD
\end{minted}

\vfill

Les \alert{transformations} (presque comme Pig).


\begin{minted}{scala}
val bergerie = loupsEtMoutons.filter(line => line.contains("bergerie"))
loupsEtMoutons.filter(line => line.contains("loup")).count() // Combinaison transf./action
\end{minted}

\end{frame}



\begin{frame}[fragile]{Le compteur de mots}

Comme promis, le compteur de mot Spark. D'abord on prend tous les termes
(\texttt{flatMap}).


\begin{minted}{scala}
val termes = loupsEtMoutons.flatMap(line => line.split(" "))
\end{minted}

\vfill

Initialisation du comptage: chaque terme compte pour 1 (\texttt{map}).


\begin{minted}{scala}
val termeUnit = termes.map(word => (word, 1))
\end{minted}

\vfill

Je regroupe les termes, et je compte (\texttt{reduce}).


\begin{minted}{scala}
val compteurTermes = termeUnit.reduceByKey((a, b) => a + b)
\end{minted}

\vfill

Et je peux rendre \alert{persistant} les compteurs de termes.


\begin{minted}{scala}
compteurTermes.persist()
\end{minted}



\end{frame}